\chapter{Commodity and Energy Derivatives}\label{ch:rc:commodity-and-energy-derivatives}

In the second week of February 2021 a cold snap froze gas wells across Texas and the
southern plains. Spot gas at the Waha hub in West Texas reached 206.19 dollars per
million British thermal units on 16 February; at one Oklahoma hub the average price on
17 February was 1\,192 dollars, against under 3 dollars in the first week of the month,
according to the US Energy Information Administration. Henry Hub, the benchmark, went
from 2.88 dollars on 1 February to 23.86 on the 17th and back under 3 by the 23rd. For a
week, the value of gas in the ground depended on how fast it could be taken out: an
owner of a salt cavern that can empty in days could sell at the peak, one with a slow
reservoir could not. This chapter values that flexibility. One Quant Book~3 described
the forward curves, storage economics and structured contracts of commodity markets;
here they get models: a forward curve that moves in two factors, options on the spread
between two prices, and the valuation of storage and swing rights by optimisation under
uncertainty.

\section{Factor models of the forward curve}

A commodity has a forward curve rather than a price: each delivery month trades, with
seasonal shape and a convenience yield (One Quant Book~3, chapter~10). A model must move
the whole curve, and it must let short-dated contracts move more than long-dated ones.

\begin{definition}[Schwartz--Smith model]\label{def:rc:commodity-and-energy-derivatives:ss}
The \emph{Schwartz--Smith model}\index{Schwartz--Smith model} writes the log spot price
as the sum of a short-term deviation $\chi_t$, an Ornstein--Uhlenbeck process reverting to
zero at speed $\kappa$ with volatility $\sigma_s$, and a long-term level $\xi_t$, a Brownian
motion with volatility $\sigma_l$, the two correlated by $\rho$. Its forward curve moves as
\[
\frac{dF(t,T)}{F(t,T)} = \sigma_se^{-\kappa(T-t)}\,dW^1_t+\sigma_l\,dW^2_t,\qquad d\langle W^1,W^2\rangle = \rho\,dt,
\]
so that $\ln F(t,T) = \ln F(0,T)+e^{-\kappa(T-t)}\chi_t+\xi_t-\tfrac12V(t,T)$ with $V$ the variance.
Seasonality sits in the initial curve $F(0,T)$, which the model fits exactly.
\end{definition}

Both factors are Gaussian and simulated exactly, and each forward is lognormal: an option
on the forward $F(\cdot,T)$ expiring at $t$ is priced by Black's formula with the implied
volatility $\sqrt{V(t,T)/t}$.

\begin{example}[Calibration]\label{ex:rc:commodity-and-energy-derivatives:calib}
An illustrative Henry Hub-like curve runs from 2.60 dollars for April to 3.50 for January.
Options expiring half a month before each delivery are quoted at implied volatilities
falling from 62.0\% (May) to 43.0\% (March), synthetic. With $\rho = 30\%$ fixed, a
Levenberg--Marquardt fit gives $\kappa = 1.50$, $\sigma_s = 60.0\%$ and $\sigma_l = 20.0\%$.
\end{example}

\section{The Samuelson effect and forward volatility}

In the model the instantaneous volatility of a forward with time $\tau$ to delivery is
\[
\sigma(\tau) = \sqrt{\sigma_s^2e^{-2\kappa\tau}+\sigma_l^2+2\rho\sigma_s\sigma_le^{-\kappa\tau}}:
\]
68.7\% at delivery, 27.2\% one year
out and 20.0\% five years out for the calibrated parameters. This is the Samuelson effect
of One Quant Book~3 as a model property: shocks to supply and demand move near contracts,
which the short-term factor carries, while far contracts move only with the long-term
level (\cref{fig:rc:commodity-and-energy-derivatives:vol}). The correlation between two
contracts falls with the distance between their deliveries, which is what makes \omterm{def:rc:commodity-and-energy-derivatives:spread}{spread
options} worth something.

\begin{omfigure}
\begin{tikzpicture}
\begin{axis}[width=11.5cm,height=5.4cm,
  xlabel={time to delivery (years)},ylabel={volatility (\%)},
  tick label style={font=\small},label style={font=\small},
  xmin=0,xmax=1,ymin=0,ymax=75,grid=major,grid style={gray!25},
  legend columns=2,legend style={font=\footnotesize,at={(0.5,-0.25)},anchor=north,draw=none,fill=none,/tikz/every even column/.append style={column sep=8pt}},
  table/col sep=comma]
\addplot[omDef,very thick] table[x=tau,y=vol]{figdata/rates-credit-risk/16-commodity-and-energy-derivatives/vol.csv};
\addplot[omThm,only marks,mark=*,mark size=1.8pt] table[x=t,y=vol]{figdata/rates-credit-risk/16-commodity-and-energy-derivatives/volquotes.csv};
\legend{instantaneous forward volatility,implied volatility quotes (at delivery date)}
\end{axis}
\end{tikzpicture}

\omcaption{The calibrated model's instantaneous forward volatility against time to
delivery, and the synthetic implied volatilities it was fitted to, plotted at each
contract's delivery time. An implied volatility averages the instantaneous volatility over
the option's life, from times to delivery of the whole horizon down to half a month, so it
lies above the curve's value at its delivery time. Data: the chapter's
tutorial.}\label{fig:rc:commodity-and-energy-derivatives:vol}
\end{omfigure}

\section{Spread options}

\begin{definition}[Spread option]\label{def:rc:commodity-and-energy-derivatives:spread}
A \emph{spread option}\index{spread option} pays $\max(F_1-F_2-K,0)$ on the difference
between two prices: two delivery months of one commodity (a calendar spread), a product
and its input (a crack or spark spread), or one commodity at two places.
\end{definition}

With $K = 0$ and lognormal prices the option is an exchange option, priced exactly by
Margrabe's formula with the volatility of the ratio,
$\sigma^2 = \sigma_1^2+\sigma_2^2-2\rho\sigma_1\sigma_2$. With $K\ne0$ the difference is not lognormal and
there is no closed form.

\begin{definition}[Kirk's approximation]\label{def:rc:commodity-and-energy-derivatives:kirk}
\emph{Kirk's approximation}\index{Kirk's approximation} prices a spread call by treating
$F_2+K$ as lognormal with volatility $\sigma_2F_2/(F_2+K)$ and applying Margrabe's formula:
\[
C \approx P\bigl(F_1N(d_1)-(F_2+K)N(d_2)\bigr),\quad \sigma_K^2 = \sigma_1^2-2\rho\sigma_1\sigma_2b+\sigma_2^2b^2,\quad b = \frac{F_2}{F_2+K}.
\]
\end{definition}

\omcode{code/firm/energymodel/firm_energymodel.py}{102}{107}{Kirk's approximation to a spread call.}\label{lst:rc:commodity-and-energy-derivatives:kirk}

\begin{example}[A calendar spread]\label{ex:rc:commodity-and-energy-derivatives:cal}
A call on the January forward (3.50) minus the July forward (2.74), strike 0.50, expiring
at July delivery: the model gives volatilities of 34.8\% and 58.1\% to that date and a
correlation of 0.970. \omterm{def:rc:commodity-and-energy-derivatives:kirk}{Kirk's approximation} prices it at 29.41 cents per MMBtu against
29.78 by Monte Carlo (standard error 0.01), an error of 0.38 cent; at zero strike both give
about 75.8 cents, Margrabe's value. The error is largest near the money, where the approximation of the
sum's distribution matters most (\cref{fig:rc:commodity-and-energy-derivatives:kirk}).
\end{example}

\begin{omfigure}
\begin{tikzpicture}
\begin{axis}[width=11.5cm,height=5.2cm,
  xlabel={strike $K$ (\$/MMBtu)},ylabel={Kirk minus Monte Carlo (cents)},
  tick label style={font=\small},label style={font=\small},
  xmin=-0.5,xmax=1.5,ymin=-0.5,ymax=0.5,grid=major,grid style={gray!25},
  legend columns=2,legend style={font=\footnotesize,at={(0.5,-0.25)},anchor=north,draw=none,fill=none,/tikz/every even column/.append style={column sep=8pt}},
  table/col sep=comma]
\addplot[omDef,very thick,mark=*,mark size=1.2pt] table[x=k,y=err]{figdata/rates-credit-risk/16-commodity-and-energy-derivatives/kirk.csv};
\addplot[omExo,dashed] table[x=k,y=se]{figdata/rates-credit-risk/16-commodity-and-energy-derivatives/kirk.csv};
\addplot[omExo,dashed,forget plot] table[x=k,y expr=-\thisrow{se}]{figdata/rates-credit-risk/16-commodity-and-energy-derivatives/kirk.csv};
\legend{Kirk's error,$\pm2$ Monte Carlo standard errors}
\end{axis}
\end{tikzpicture}

\omcaption{Error of \omterm{def:rc:commodity-and-energy-derivatives:kirk}{Kirk's approximation} for the January--July calendar spread call
against 400\,000 Monte Carlo paths, by strike. The error is zero at zero strike (Margrabe
is exact) and changes sign across the money; it stays under half a cent at every strike,
on options worth from nothing to 1.26 dollars. Data: the chapter's tutorial.}\label{fig:rc:commodity-and-energy-derivatives:kirk}
\end{omfigure}

\section{Storage valuation}

A storage facility buys gas when it is cheap, holds it and sells it when it is dear,
within limits on capacity and on injection and withdrawal rates, and at a cost per unit
moved. Its owner holds a portfolio of calendar \omterm{def:rc:commodity-and-energy-derivatives:spread}{spread options} linked by the inventory.

\begin{definition}[Intrinsic storage value]\label{def:rc:commodity-and-energy-derivatives:intrinsic}
The \emph{intrinsic storage value}\index{intrinsic storage value} of a facility is the
value of the best injection and withdrawal schedule on today's forward curve, locked in
by trading the forwards: a deterministic optimisation.
\end{definition}

\begin{definition}[Extrinsic value, rolling intrinsic]\label{def:rc:commodity-and-energy-derivatives:extrinsic}
The \emph{extrinsic value}\index{extrinsic value} of a facility is its value above the
intrinsic, from the right to change the schedule as prices move. The \emph{rolling
intrinsic}\index{rolling intrinsic} strategy captures part of it: re-optimise the schedule on
each new forward curve and re-hedge the difference, never losing the intrinsic already
locked in.
\end{definition}

\begin{method}[Storage by least-squares Monte Carlo]\label{met:rc:commodity-and-energy-derivatives:lsm}
Simulate spot prices for each period. Work backwards from the last period with the value
of each end inventory known on each path. At each period and for each inventory level
reachable after the action, regress the next period's realised value on $1$, $S$ and $S^2$
to estimate the continuation value; choose, path by path, the action that maximises cash
now plus estimated continuation, and carry back the realised value of that choice. The
value is the average at the starting inventory.
\end{method}

\omcode{code/firm/energymodel/firm_energymodel.py}{169}{200}{Least-squares Monte Carlo for storage: regression of continuation values per end inventory, and realised values carried back.}\label{lst:rc:commodity-and-energy-derivatives:lsm}

\begin{example}[A season of storage]\label{ex:rc:commodity-and-energy-derivatives:storage}
A facility of 1 million MMBtu (ten units) starts and ends the April--March year empty; it
can inject two units a month and withdraw four, at 2 cents per MMBtu each way. On the
curve of \cref{ex:rc:commodity-and-energy-derivatives:calib} the intrinsic plan injects
two units a month from April to August and withdraws in December, January and February
(\cref{fig:rc:commodity-and-energy-derivatives:curve}); its value is USD~634\,447. \omterm{def:rc:commodity-and-energy-derivatives:extrinsic}{Rolling
intrinsic} on 4\,000 simulated curves is worth USD~699\,550, least-squares Monte Carlo
USD~698\,182: an \omterm{def:rc:commodity-and-energy-derivatives:extrinsic}{extrinsic value} of about 10\% of the intrinsic at a monthly decision
frequency. Daily decisions, which a fast facility can take, add more.
\end{example}

\begin{omfigure}
\begin{tikzpicture}
\begin{axis}[width=11.5cm,height=5.4cm,
  xlabel={delivery month},ylabel={forward (\$/MMBtu)},
  axis y line*=left,
  xtick={1,2,3,4,5,6,7,8,9,10,11,12},xticklabels={Apr,May,Jun,Jul,Aug,Sep,Oct,Nov,Dec,Jan,Feb,Mar},
  tick label style={font=\small},label style={font=\small},
  xmin=0.5,xmax=12.5,ymin=2.4,ymax=3.6,
  table/col sep=comma]
\addplot[omDef,very thick,mark=*,mark size=1.5pt] table[x=month,y=forward]{figdata/rates-credit-risk/16-commodity-and-energy-derivatives/curve.csv};\label{pl:rc:energy:fwd}
\end{axis}
\begin{axis}[width=11.5cm,height=5.4cm,
  axis y line*=right,axis x line=none,ylabel={inventory (units)},
  tick label style={font=\small},label style={font=\small},
  xmin=0.5,xmax=12.5,ymin=0,ymax=12,
  legend columns=2,legend style={font=\footnotesize,at={(0.5,-0.25)},anchor=north,draw=none,fill=none,/tikz/every even column/.append style={column sep=8pt}},
  table/col sep=comma]
\addlegendimage{/pgfplots/refstyle=pl:rc:energy:fwd}\addlegendentry{forward curve}
\addplot[omThm,very thick,const plot mark mid] table[x=month,y=inventory]{figdata/rates-credit-risk/16-commodity-and-energy-derivatives/curve.csv};
\addlegendentry{intrinsic inventory at month end}
\end{axis}
\end{tikzpicture}

\omcaption{The illustrative forward curve (left axis) and the intrinsic plan's inventory
at the end of each month (right axis): fill in the cheap summer months, empty in the dear
winter ones. Data: the chapter's tutorial.}\label{fig:rc:commodity-and-energy-derivatives:curve}
\end{omfigure}

\section{Swing contracts}

A swing contract (One Quant Book~3, chapter~12) is storage without injection: the holder
has a number of rights to take gas at a fixed price, limited per period and in total. The
same dynamic programme values it, with the inventory counting rights left.

\begin{example}[A winter swing]\label{ex:rc:commodity-and-energy-derivatives:swing}
Rights to take up to 200\,000 MMBtu a month from November to March, 600\,000 in all, at
3.10 dollars: on the forward curve the best plan takes 200\,000 in December, January and
February, an intrinsic value of 1.84 per unit of 100\,000 MMBtu; least-squares Monte Carlo
gives 4.37. Ten monthly calls, two a month without the total limit, are worth 6.11, an
upper bound: the limit on the total is what makes the swing cheaper than a strip.
\end{example}

\begin{dated}{2026-09}{February 2021}\label{dat:rc:commodity-and-energy-derivatives:uri}
The US Energy Information Administration reported that during the cold snap of February
2021 the Waha hub reached 206.19 dollars per MMBtu on 16 February and SoCal Citygate 144.00
on 12 February, and that Oneok Gas Transportation in Oklahoma averaged 1\,192 dollars on
17 February according to NGI data, against 2.91 in the first week of the month; gas
production fell because of well freeze-offs. Henry Hub averaged 5.35 dollars over the month, against
2.71 in January.
\end{dated}

\begin{omfigure}
\begin{tikzpicture}
\begin{axis}[width=11.5cm,height=5.2cm,
  ylabel={Henry Hub spot (\$/MMBtu)},
  xlabel={December 2020 to April 2021},
  xtick={-31,0,31,59,90,120},xticklabels={Dec,Jan,Feb,Mar,Apr,May},
  tick label style={font=\small},label style={font=\small},
  xmin=-31,xmax=120,ymin=0,ymax=26,grid=major,grid style={gray!25},
  table/col sep=comma]
\addplot[omDef,thick,mark=*,mark size=0.8pt] table[x=day,y=price]{figdata/rates-credit-risk/16-commodity-and-energy-derivatives/henryhub.csv};
\node[omThm,font=\footnotesize,anchor=west] at (axis cs:49,23.86) {23.86 on 17 February};
\end{axis}
\end{tikzpicture}

\omcaption{Henry Hub daily spot price, December 2020 to April 2021. The cold snap lifted
the price from under 3 dollars to 23.86 for one day and it was back under 3 within a week:
value that only a facility able to deliver within days could capture. Data: US Energy
Information Administration, series RNGWHHD (public domain).}\label{fig:rc:commodity-and-energy-derivatives:hh}
\end{omfigure}

\section{Tutorial: storage under a two-factor curve}\label{tut:rc:commodity-and-energy-derivatives:tutorial}

\begin{tutorial}
\textbf{Goal.} Calibrate the two-factor model, price a calendar \omterm{def:rc:commodity-and-energy-derivatives:spread}{spread option}, and value a
storage facility and a swing contract. \textbf{End state:} the numbers of
\cref{ex:rc:commodity-and-energy-derivatives:calib,ex:rc:commodity-and-energy-derivatives:cal,ex:rc:commodity-and-energy-derivatives:storage,ex:rc:commodity-and-energy-derivatives:swing}
and the four figures.
\begin{tutsteps}
\item \textbf{Calibrate}: \texttt{model()} fits $\kappa$, $\sigma_s$, $\sigma_l$ to
\texttt{vol\_quotes()}.
\item \textbf{Spread}: \texttt{calendar\_spread(0.5)}; \texttt{kirk\_table()} by strike.
\item \textbf{Storage}: \texttt{storage()} gives intrinsic, \omterm{def:rc:commodity-and-energy-derivatives:extrinsic}{rolling intrinsic} and
least-squares Monte Carlo values.
\item \textbf{Swing and storm}: \texttt{swing()}, \texttt{storm()};
\texttt{fig\_rc\_energy.py} writes the charts.
\end{tutsteps}
\textbf{What to change next.} Double the withdrawal rate and see how much \omterm{def:rc:commodity-and-energy-derivatives:extrinsic}{extrinsic value}
the faster facility adds; set $\sigma_s$ to zero and watch the \omterm{def:rc:commodity-and-energy-derivatives:extrinsic}{extrinsic value} of the
storage shrink.
\end{tutorial}

\section{Build: the energy model}\label{bld:rc:commodity-and-energy-derivatives:energymodel}

\begin{build}
\bfield{Purpose} The firm's commodity engine: forward-curve simulation for the exposure and
risk engines of Parts~III and~IV, \omterm{def:rc:commodity-and-energy-derivatives:spread}{spread options}, and the valuation of physical
flexibility.
\bfield{Interface} \texttt{TwoFactor(kappa, sigma\_s, sigma\_l, rho)} with \texttt{var},
\texttt{cov}, \texttt{implied\_vol}, \texttt{simulate}, \texttt{forward};
\texttt{calibrate}; \texttt{margrabe}, \texttt{kirk}, \texttt{spread\_mc};
\texttt{Facility}; \texttt{intrinsic}, \texttt{lsm}, \texttt{rolling\_intrinsic},
\texttt{spot\_paths}; \texttt{swing\_facility}.
\bfield{Rules} One decision per delivery period; integer inventory units; spot of a period
is its forward at delivery; antithetic simulation.
\bfield{Acceptance tests} \texttt{code/firm/energymodel/tests/}: simulated forwards are
martingales with the model's variance; calibration recovers the parameters; Kirk equals
Margrabe at zero strike and is within 1\% of Monte Carlo; a flat curve has no \omterm{def:rc:commodity-and-energy-derivatives:intrinsic}{intrinsic
storage value}; without volatility, least-squares Monte Carlo equals intrinsic; \omterm{def:rc:commodity-and-energy-derivatives:extrinsic}{rolling
intrinsic} and least-squares Monte Carlo exceed intrinsic; a swing's intrinsic is the sum of
its best exercises.
\bfield{Stretch} Daily decisions within the month; ratchets (rates that depend on the
inventory); a jump factor for spikes; the spot-forward basis of a physical hub.
\end{build}

\begin{omsources}
\item E.\ S.\ Schwartz, ``The stochastic behavior of commodity prices: implications for
valuation and hedging'', \emph{Journal of Finance} 52(3), 1997, 923--973.
\item E.\ Schwartz and J.\ E.\ Smith, ``Short-term variations and long-term dynamics in
commodity prices'', \emph{Management Science} 46(7), 2000, 893--911.
\item A.\ Boogert and C.\ de Jong, ``Gas storage valuation using a Monte Carlo method'',
\emph{Journal of Derivatives} 15(3), 2008, 81--98.
\item E.\ Kirk, ``Correlation in the energy markets'', in R.\ Jameson (ed.), \emph{Managing Energy Price
Risk}, Risk Publications and Enron, London, 1995, 71--78.
\item US Energy Information Administration, \emph{Today in Energy}, ``Cold weather brings
near record-high natural gas spot prices'', February 2021; Henry Hub daily spot series.
\end{omsources}

\section{Exercises}

\begin{exercise}[$\star$]\label{exo:rc:commodity-and-energy-derivatives:1}
With the calibrated parameters, what is the instantaneous volatility of a forward one year
from delivery, and why is it lower than at delivery?
\end{exercise}

\begin{exercise}[$\star$]\label{exo:rc:commodity-and-energy-derivatives:2}
Why is the intrinsic value of storage zero on a flat forward curve, and why is the
facility still worth something?
\end{exercise}

\begin{exercise}[$\star$]\label{exo:rc:commodity-and-energy-derivatives:3}
A calendar spread call with strike 0.50 on forwards at 3.50 and 2.74 expires tomorrow.
What is it worth?
\end{exercise}

\begin{exercise}[$\star\star$]\label{exo:rc:commodity-and-energy-derivatives:4}
Why does the correlation between two delivery months matter for a \omterm{def:rc:commodity-and-energy-derivatives:spread}{spread option}, and in
which direction?
\end{exercise}

\begin{exercise}[$\star\star$]\label{exo:rc:commodity-and-energy-derivatives:5}
Why can \omterm{def:rc:commodity-and-energy-derivatives:extrinsic}{rolling intrinsic} never be worth less than intrinsic?
\end{exercise}

\begin{exercise}[$\star\star$]\label{exo:rc:commodity-and-energy-derivatives:6}
Why is a swing contract worth less than the strip of monthly calls it contains?
\end{exercise}

\begin{exercise}[$\star\star\star$]\label{exo:rc:commodity-and-energy-derivatives:7}
\emph{Coding.} Compute Kirk's error at strikes 0 and 0.5 and explain the difference.
\end{exercise}

\begin{exercise}[$\star\star\star$]\label{exo:rc:commodity-and-energy-derivatives:8}
\emph{Find the flaw.} ``Our storage book is valued at intrinsic and fully hedged with
forwards, so February 2021 was neither a risk nor an opportunity for us.''
\end{exercise}

\section{Problem: The Salt Cavern in February 2021}

\begin{problem}[{Weekend problem --- a fast and a slow facility in the cold snap}]\label{pb:rc:commodity-and-energy-derivatives:1}
Two facilities each hold 500\,000 MMBtu near Henry Hub on 8 February 2021. The salt cavern
can withdraw 100\,000 MMBtu a day, the depleted reservoir 10\,000. Between 8 and 19 February
each may sell on its best days at the daily Henry Hub price, and buys the same volume back
at the March 2021 average to restore its plan; moving gas costs 2 cents per MMBtu each way.

\medskip\noindent\textbf{Part I --- Before the winter.}
\begin{enumerate}
\item Give the intrinsic, \omterm{def:rc:commodity-and-energy-derivatives:extrinsic}{rolling intrinsic} and least-squares Monte Carlo values of the
chapter's facility.
\item What is its \omterm{def:rc:commodity-and-energy-derivatives:extrinsic}{extrinsic value}, and where does it come from?
\item Which parameter of the curve model drives the \omterm{def:rc:commodity-and-energy-derivatives:extrinsic}{extrinsic value} most?
\item Why does a monthly model understate a salt cavern's \omterm{def:rc:commodity-and-energy-derivatives:extrinsic}{extrinsic value}?
\item How would you hedge the intrinsic value at the start of the season?
\end{enumerate}

\medskip\noindent\textbf{Part II --- The storm.}
\begin{enumerate}[resume]
\item Give Henry Hub on 1 February and at its peak, and the day of the peak.
\item Give the March 2021 average used for the buy-back.
\item Give the volume the salt cavern sells, on how many days, and its gain.
\item Give the same for the depleted reservoir.
\item Why is the ratio of the gains not the ratio of the withdrawal rates?
\end{enumerate}

\medskip\noindent\textbf{Part III --- What a model sees.}
\begin{enumerate}[resume]
\item How many standard deviations is the move from 3.76 dollars on 10 February to 23.86
on 17 February (four trading days) under a lognormal model with the chapter's 69\%
short-term volatility?
\item Which model feature would put such spikes in the simulation?
\item Why did the storm raise the value of the facility's options but not its intrinsic
value for the next season?
\item Which contracts gained besides storage?
\item Who was short that flexibility?
\end{enumerate}

\medskip\noindent\textbf{Part IV --- Judgement.}
\begin{enumerate}[resume]
\item Why do owners of fast storage pay for it when the average year gives little
\omterm{def:rc:commodity-and-energy-derivatives:extrinsic}{extrinsic value}?
\item How would you report the storage book's risk to a risk committee?
\item What would you change in the valuation model after February 2021?
\item State the \emph{named result}: the gains of the fast and the slow facility in the
storm.
\item In one sentence: what is a storage facility worth beyond its seasonal spread?
\end{enumerate}
\end{problem}

\section{Interview questions}

\begin{interviewq}[$\star$ \iqroles{trader, researcher}]\label{iq:rc:commodity-and-energy-derivatives:1}
What is the Samuelson effect, and how does a two-factor model produce it?
\end{interviewq}

\begin{interviewq}[$\star\star$ \iqroles{researcher}]\label{iq:rc:commodity-and-energy-derivatives:2}
Price a \omterm{def:rc:commodity-and-energy-derivatives:spread}{spread option} with a non-zero strike. What approximations exist and when do they
fail?
\end{interviewq}

\begin{interviewq}[$\star\star$ \iqroles{trader}]\label{iq:rc:commodity-and-energy-derivatives:3}
How do you value a gas storage facility? Distinguish intrinsic, \omterm{def:rc:commodity-and-energy-derivatives:extrinsic}{rolling intrinsic} and full
optionality.
\end{interviewq}

\begin{interviewq}[$\star\star$ \iqroles{developer, researcher}]\label{iq:rc:commodity-and-energy-derivatives:4}
Describe least-squares Monte Carlo for storage. What are its biases?
\end{interviewq}

\begin{interviewq}[$\star\star\star$ \iqroles{risk}]\label{iq:rc:commodity-and-energy-derivatives:5}
How would you risk-manage a book of storage and swing contracts through a winter?
\end{interviewq}

\begin{interviewq}[$\star\star\star$ \iqroles{researcher}]\label{iq:rc:commodity-and-energy-derivatives:6}
Why does a Gaussian two-factor model miss power and gas price spikes, and what would you
add?
\end{interviewq}
